\section{Vai trò của RAG trong hệ thống tư vấn tuyển sinh}

Phân hệ RAG, viết tắt của Retrieval-Augmented Generation, là thành phần chịu trách nhiệm xử lý các câu hỏi cần tri thức từ tài liệu tuyển sinh, chương trình đào tạo, điểm chuẩn, quy chế, biểu mẫu và các văn bản liên quan của Học viện. Khác với Rasa, vốn phù hợp với các intent định sẵn và các luồng hội thoại có cấu trúc, RAG được dùng khi câu trả lời cần dựa trên kho tài liệu đã được ingest, chunking và lập chỉ mục.

Trong hệ thống KMA Chatbot, RAG không được thiết kế như một mô hình tự trả lời theo trí nhớ của LLM. Thay vào đó, RAG hoạt động theo nguyên tắc: trước khi sinh câu trả lời, hệ thống phải truy hồi bằng chứng từ kho tài liệu, xây dựng ngữ cảnh, sau đó mới gọi LLM để diễn đạt câu trả lời. Cách thiết kế này giúp giảm nguy cơ hallucination, tăng khả năng truy vết nguồn và phù hợp với bài toán tuyển sinh, nơi thông tin như phương thức xét tuyển, điểm chuẩn, học phí hoặc chương trình đào tạo có thể thay đổi theo từng năm.

RAG không để LLM tự trả lời theo tri thức đã học sẵn, mà trước hết truy xuất thông tin liên quan từ kho tài liệu, sau đó đưa các bằng chứng này vào mô hình sinh để tạo câu trả lời có căn cứ.

\begin{figure}[H]
    \centering
    \includegraphics[width=0.9\textwidth]{\detokenize{content/source/RAG/mô hình RAG.png}}
    \caption{Kiến trúc tổng quan của mô hình RAG}
    \label{fig:rag-basic-model}
\end{figure}

Trong dữ liệu kỹ thuật đã khảo sát, hệ thống có frontend web, backend API, RAG service ở port 9621 và các lớp lưu trữ phía sau. RAG service phụ trách ingest, chunking, indexing, retrieval và sinh câu trả lời theo tài liệu. Health endpoint cho thấy phân hệ RAG sử dụng embedding model BAAI/bge-m3, RedisKVStorage, RedisDocStatusStorage, MemgraphStorage và QdrantVectorDBStorage. Cấu trúc này cho thấy metadata, trạng thái tài liệu, vector và đồ thị tri thức được tách thành các lớp lưu trữ riêng.

Về mặt chức năng, RAG đảm nhiệm các nhóm nhiệm vụ chính sau:

\begin{itemize}
    \item Tiếp nhận và xử lý tài liệu tuyển sinh, đào tạo, điểm chuẩn, biểu mẫu.
    \item Chia tài liệu thành các chunk có thể truy hồi.
    \item Sinh embedding cho chunk, entity và relation.
    \item Lưu trữ vector vào vector database.
    \item Lưu trạng thái xử lý tài liệu để quản trị viên theo dõi.
    \item Truy hồi bằng chứng khi người dùng đặt câu hỏi.
    \item Xây dựng context và gọi LLM sinh câu trả lời.
    \item Trả lời kèm nguồn hoặc căn cứ tài liệu.
\end{itemize}

Như vậy, RAG là phân hệ bổ sung tri thức động cho chatbot. Nếu Rasa giúp hệ thống hiểu và điều phối hội thoại theo intent, thì RAG giúp hệ thống trả lời các câu hỏi có nội dung phụ thuộc tài liệu thực tế.
